\documentclass{article}
\usepackage{graphicx} % Required for inserting images
  \RequirePackage[french]{babel}
  \RequirePackage[T1]{fontenc}
  \RequirePackage[utf8]{inputenc}
  
\title{STT-2902 - Modélisation statistique}
\author{Julien Miron}


\begin{document}

\section*{Équipe 1 – Statistiques descriptives et mauvaises interprétations}

Objectif : Montrer comment des statistiques correctement calculées peuvent être mal interprétées ou manipulées. \textit{Notez que vos collègues ne savent pas que vous devrez les induire en erreur, vous pouvez utiliser ceci pour causer un effet de «surprise».}

\vspace{1cm}
Suggestions :
\begin{itemize}
    \item Préparez un quiz avec pièges : présentez des graphiques ou des moyennes trompeuses et demandez aux autres équipes ce qu’elles concluraient.
    \item Incluez de vrais exemples tirés de médias ou de publicités.

\end{itemize}

\vspace{1cm}
\textbf{Requis} :
\begin{enumerate}
    \item Inclure au moins un graphique où l’échelle ou la représentation visuelle pourrait induire en erreur.
    \item Utiliser la moyenne dans un contexte où la médiane serait un meilleur indicateur.
\end{enumerate}


\end{document}
